\section{LOCAL\_MODE: arquitectura local}
\label{sec:local_mode}

LOCAL\_MODE está pensado para operar sobre una banda transportadora
física, en un banco de pruebas sin acceso a internet.

\subsection{Topología de red}
El propio teléfono Android que toma las fotografías crea una red WiFi
(hotspot); no existe un enrutador independiente. Tanto el
microcontrolador ESP32 como la computadora que ejecuta el modelo de
clasificación se conectan a esa red como clientes ordinarios. Al no
existir mecanismo de reserva de IP fija como en un enrutador
convencional, la dirección de la computadora se configura de forma
manual (IP estática, o leyendo la concesión DHCP desde la
configuración del hotspot). El ESP32, por restricción de hardware,
solo opera en banda de 2.4~GHz, por lo que el hotspot debe forzarse a
ese modo.

\subsection{Control de la banda transportadora}
El movimiento de la banda está controlado por un ESP32-WROOM-32 que
acciona un puente H (BTS7960). El motor solo se habilita cuando un
sensor de huella dactilar reconoce una huella previamente registrada:
cada huella autorizada alterna el estado de la banda (encendida o
apagada), de forma que el sistema nunca mueve la banda para un usuario
no autorizado. El estado del sistema (banda encendida/apagada, última
clasificación, diámetro estimado) se muestra en una pantalla LCD de
16$\times$2 caracteres.

\subsection{Detección de pieza y captura fotográfica}
Mientras la banda está en movimiento, un sensor detecta cuándo una
pieza se encuentra bajo la cámara y detiene momentáneamente el
movimiento, dando tiempo a que el teléfono capture la fotografía antes
de reanudar el recorrido. Esta pausa es gestionada enteramente por el
sistema, de forma automática, sin intervención manual del operario.

\subsection{Servidor local y clasificación}
Un servidor FastAPI (\textit{local-server}), sin autenticación ni TLS
por tratarse de una red cerrada, expone un pequeño conjunto de
endpoints HTTP: recepción de la fotografía capturada, consulta del
resultado más reciente (que es lo que el ESP32 consulta
periódicamente, al no existir un mecanismo de notificación push en
este modo) y verificación de estado del servicio.

Sobre cada fotografía recibida se ejecutan dos procesos
independientes:

\begin{itemize}
  \item \textbf{Clasificación}: un modelo YOLOv8 entrenado
  específicamente sobre fotografías de las piezas del banco de pruebas
  (brocas, escariadores, tornillos y brocas de centropunto) determina
  la clase de la pieza y su nivel de confianza.
  \item \textbf{Medición de diámetro}: a partir de un marcador de
  referencia de longitud física conocida (1~cm) colocado junto a la
  pieza en el mismo encuadre, se calcula la relación
  píxeles-por-centímetro y, con ella, el diámetro real de la pieza
  fotografiada.
\end{itemize}

Ambos resultados se combinan en una única respuesta, que el ESP32
consulta y muestra en su pantalla LCD y por puerto serial.
